\documentclass[11pt]{article}

\usepackage[margin=1in]{geometry}
\usepackage{amsmath,amssymb}
\usepackage{booktabs}
\usepackage{enumitem}
\usepackage{microtype}
\usepackage[hidelinks]{hyperref}

\setlength{\parindent}{0pt}
\setlength{\parskip}{0.7em}
\setlist{itemsep=0.3em,topsep=0.4em}

\title{Metropolis--Hastings:\\A Multidisciplinary Story of Modeling and Simulation}
\author{PHYS 441}
\date{}

\begin{document}

\maketitle

\begin{abstract}
The Metropolis--Hastings algorithm is now used in physics, statistics, computer
science, engineering, and many other fields. Its origin is an unusually clear
example of how modeling and simulation develop: a physical question motivated a
mathematical reformulation; an algorithm made that reformulation practical; a
program translated the algorithm into machine instructions; and a new computer
made the calculation possible. This reading follows that development from Los
Alamos in the 1940s and 1950s to Hastings's generalization in 1970. It also
examines a historical complication: the familiar name of the algorithm does not
fully describe the contributions of the people who created it.
\end{abstract}

\section{Why this history belongs in a modeling course}

A finished simulation can make its component parts look inevitable. We are
given a model, write an algorithm, run a program, and interpret the output. The
history of Metropolis--Hastings shows that these steps are neither automatic nor
the work of a single discipline.

The original project required all of the following:

\begin{itemize}
  \item a \textbf{physical model} of matter in thermal equilibrium;
  \item a \textbf{statistical description} of which configurations should be
        common and which should be rare;
  \item a \textbf{mathematical insight} that the desired average could be found
        without evaluating an enormous integral directly;
  \item an \textbf{algorithm} that generated configurations with the correct
        long-run frequencies;
  \item a \textbf{program} written for a machine with extremely limited memory;
  \item a functioning \textbf{electronic computer}; and
  \item scientific judgment about validation, approximation, and uncertainty.
\end{itemize}

Physics, computer science, mathematics, and engineering were not separate
layers added after the fact. The method arose from their interaction.

\section{The Los Alamos setting}

Modern Monte Carlo computation grew out of the scientific environment created
by the Manhattan Project. Many Los Alamos calculations involved systems with a
vast number of possible histories or configurations: neutron transport,
radiation transport, thermonuclear reactions, hydrodynamics, and equations of
state. Exact analytic solutions were rarely available, while exhaustive
numerical enumeration was impossible.

During the war, much numerical work was performed by teams of human
``computers'' and electromechanical machines. Electronic machines such as ENIAC
then demonstrated that calculations could be performed at an unprecedented
scale. Stanis\l aw Ulam recognized that this new speed could revive statistical
sampling methods that had previously required too much arithmetic. John von
Neumann helped formulate and implement early applications, including neutron
transport calculations \cite{metropolis1987}.

The name \emph{Monte Carlo} emphasized the use of random numbers and alluded to
games of chance. Randomness, however, was not being used as a substitute for
rigor. It was being engineered into a numerical method: repeat a random
experiment, calculate an observable, and use the resulting statistics to
estimate a quantity that is otherwise difficult to compute.

\section{The people behind the 1953 paper}

The landmark paper, ``Equation of State Calculations by Fast Computing
Machines,'' was published in 1953 by Nicholas Metropolis, Arianna W.
Rosenbluth, Marshall N. Rosenbluth, Augusta H. Teller, and Edward Teller
\cite{metropolis1953}. The author list brought together two married couples and
the leader of the Los Alamos computing effort.

\subsection{Nicholas Metropolis: physics and computing infrastructure}

Nicholas Constantine Metropolis (1915--1999) trained as an experimental
physicist at the University of Chicago. He worked with Enrico Fermi, joined the
Manhattan Project at Los Alamos in 1943, and worked on equations of state for
materials at high temperature and pressure. After the war, he became a major
figure in scientific computing.

Metropolis directed the construction of the Los Alamos MANIAC computer, a
stored-program machine based on the architecture developed by von Neumann and
his collaborators. MANIAC began operating in 1952 and performed the calculation
reported in the 1953 paper. Metropolis supplied far more than a piece of
hardware: he led the group, secured institutional support, allocated scarce
machine time, and helped create a setting in which physicists and programmers
could turn new numerical ideas into working calculations \cite{anderson1986}.

\subsection{Marshall Rosenbluth: the statistical-mechanical algorithm}

Marshall Rosenbluth (1927--2003) was a theoretical physicist who later became a
leading figure in plasma physics and controlled-fusion research. In his later
account of the 1953 project, he credited Edward Teller with suggesting the use
of Monte Carlo methods to take statistical-mechanical ensemble averages.
Rosenbluth then worked out the acceptance procedure that made this idea
efficient \cite{gubernatis2005}.

The essential insight was to avoid drawing configurations uniformly and then
assigning almost all of them negligible statistical weight. Instead, the
algorithm would construct a random walk that spent more time in configurations
with larger Boltzmann weight.

\subsection{Arianna Rosenbluth: turning the method into a program}

Arianna Wright Rosenbluth (1927--2020) was a physicist and programmer. She
wrote the MANIAC implementation used for the published calculation. Programming
an early computer was not a matter of translating a few lines of pseudocode into
a high-level language. The machine had limited memory, primitive input and
output, and no modern operating system or numerical library. The programmer had
to understand the mathematics, the machine representation, the control flow,
and the numerical limitations together.

Her work is a reminder that an algorithm is not yet a simulation. A simulation
exists only after the algorithm has been expressed in a form that a real machine
can execute and its results have been checked. Arianna Rosenbluth's role was
underemphasized for decades, in part because programming was often treated as
secondary labor and because women in early computing were frequently omitted
from later accounts \cite{apsrosenbluth}.

\subsection{Augusta and Edward Teller: computation and physical motivation}

Augusta ``Mici'' Teller (1909--2000) studied mathematics in Budapest and worked
in the Los Alamos computations group during the Manhattan Project. She later
worked with early computers at Argonne and Los Alamos. Accounts suggest that
she participated in the initial programming effort, although Arianna Rosenbluth
reported that she subsequently rewrote the program used for the calculation
\cite{gubernatis2005,micibio}.

Edward Teller (1908--2003) was a theoretical physicist deeply involved in
nuclear and thermonuclear research. His interest in the project was scientifically
natural. Statistical mechanics relates microscopic particle configurations to
macroscopic properties such as pressure, temperature, and density. Equations of
state are also essential in modeling matter under the extreme conditions found
in shocks, dense plasmas, stellar interiors, and nuclear-weapon calculations.

The published demonstration was not itself a simulation of a nuclear weapon.
It studied a simplified equilibrium system of interacting particles. It did,
however, emerge from a laboratory in which equations of state, high-dimensional
calculations, and the development of nuclear technology were closely connected.

\section{The physical and mathematical problem}

Consider a system of $N$ interacting particles. A single state $x$ specifies
the positions of all particles, and $E(x)$ is the energy of that configuration.
In the canonical ensemble, a configuration has probability density proportional
to its Boltzmann factor,

\begin{equation}
  \widetilde{\pi}(x)=\exp\!\left[-\frac{E(x)}{k_{\mathrm B}T}\right],
\end{equation}

where $T$ is temperature and $k_{\mathrm B}$ is Boltzmann's constant. The
equilibrium expectation of an observable $F$ is

\begin{equation}
 \langle F\rangle =
 \frac{\displaystyle\int F(x)e^{-E(x)/(k_{\mathrm B}T)}\,dx}
      {\displaystyle\int e^{-E(x)/(k_{\mathrm B}T)}\,dx}.
 \label{eq:ensemble-average}
\end{equation}

If each of hundreds of particles contributes two or three coordinates, then
the integral is hundreds-dimensional. A grid with only ten values per coordinate
would require $10^{200}$ points for a 200-dimensional problem. This is the
\emph{curse of dimensionality}: increasing computational power alone cannot make
such a grid practical.

Naive random sampling is also inefficient. Most randomly selected configurations
may have particles unrealistically close together, producing very large energies
and nearly zero Boltzmann weights. The computer spends almost all of its effort
generating configurations that contribute almost nothing to
Eq.~\eqref{eq:ensemble-average}.

\section{The algorithmic leap}

The 1953 method generates a sequence of configurations

\[
 X_0,X_1,X_2,\ldots
\]

instead of selecting independent configurations uniformly. Starting from the
current configuration $x$, it performs the following steps:

\begin{enumerate}
  \item Select a particle and propose a small random displacement, producing a
        candidate configuration $y$.
  \item Calculate the energy change $\Delta E=E(y)-E(x)$.
  \item If $\Delta E\leq 0$, accept the move.
  \item If $\Delta E>0$, accept the move with probability
        $e^{-\Delta E/(k_{\mathrm B}T)}$.
  \item If the proposal is rejected, record the old configuration again.
\end{enumerate}

For a symmetric proposal, this is equivalent to accepting with probability

\begin{equation}
 \alpha(x,y)=\min\!\left\{1,
 \frac{\widetilde{\pi}(y)}{\widetilde{\pi}(x)}\right\}.
 \label{eq:metropolis}
\end{equation}

The unknown normalization integral cancels in the ratio. More importantly, the
accepted configurations occur with the desired equilibrium frequencies. After
initialization effects have diminished, the integral in
Eq.~\eqref{eq:ensemble-average} can be estimated by the ordinary average

\begin{equation}
 \langle F\rangle \approx \frac{1}{M}\sum_{t=1}^{M}F(X_t).
\end{equation}

This is a profound change in the computational problem. Rather than evaluating
an impossible integral, the machine performs local energy calculations and
random accept--reject decisions. The answer emerges from the long-run behavior
of the constructed Markov chain.

The original paper demonstrated the method using two-dimensional hard disks.
The model was deliberately simple: disks could not overlap, and the calculation
estimated an equation of state. Simplicity made the method testable while still
addressing the collective behavior of many interacting particles.

\section{From Metropolis to Metropolis--Hastings}

The proposal in the 1953 calculation was symmetric. If the mechanism could
propose $y$ from $x$, the reverse displacement was equally likely. In 1970,
W. Keith Hastings generalized the method to an arbitrary proposal density
$q(y\mid x)$ \cite{hastings1970}. The acceptance probability became

\begin{equation}
 \alpha(x,y)=\min\!\left\{1,
 \frac{\widetilde{\pi}(y)q(x\mid y)}
      {\widetilde{\pi}(x)q(y\mid x)}\right\}.
 \label{eq:hastings}
\end{equation}

The extra ratio corrects for a proposal mechanism that favors one direction.
When the proposal is symmetric, $q(y\mid x)=q(x\mid y)$, the proposal terms
cancel, and Eq.~\eqref{eq:hastings} reduces to the original Metropolis rule in
Eq.~\eqref{eq:metropolis}.

Hastings transformed a powerful statistical-physics procedure into a broadly
applicable sampling framework. The method now supports Bayesian inference,
molecular simulation, lattice field theory, reliability analysis, inverse
problems, uncertainty quantification, and many other applications.

\section{Whose algorithm is it?}

The familiar name creates a misleadingly simple story. Historical recollections
do not agree in every detail. Edward Teller later described a strongly
collaborative effort, while Marshall Rosenbluth assigned the central algorithmic
and programming work primarily to himself and Arianna. Later historical work
generally supports the conclusion that the Rosenbluths made contributions much
larger than the usual name suggests \cite{gubernatis2005}.

Why, then, did ``Metropolis algorithm'' become standard? Metropolis was first
author, led the computing group, and made MANIAC available. Naming a result after
the first author was natural, and a short name spread more easily than the full
author list. The episode illustrates several features of scientific credit:

\begin{itemize}
  \item the name attached to a method need not identify its sole inventor;
  \item programming and implementation may be undervalued relative to theory;
  \item institutional leadership can be essential even when it is not the same
        as inventing the central formula; and
  \item later recollections must be compared critically rather than treated as
        perfectly objective records.
\end{itemize}

The aim is not to replace one heroic inventor with another. The better lesson is
that the method was produced by a network of complementary contributions.

\section{A timeline of the development}

\begin{center}
\begin{tabular}{@{}p{0.13\textwidth}p{0.78\textwidth}@{}}
\toprule
\textbf{Period} & \textbf{Development} \\
\midrule
1940s & Los Alamos uses large-scale numerical calculation for neutron transport,
        weapons physics, and equations of state. Ulam and von Neumann develop
        modern Monte Carlo ideas for electronic computers. \\
1948--1952 & Metropolis leads the construction of MANIAC at Los Alamos. \\
1952--1953 & Teller, the Rosenbluths, the Tellers, and Metropolis develop and run
             the statistical-mechanics calculation on MANIAC. \\
1953 & The five-author paper introduces the acceptance algorithm using a
       two-dimensional hard-disk equation-of-state calculation. \\
1970 & Hastings publishes the general acceptance rule for asymmetric proposals. \\
Later decades & MCMC becomes central to computational physics and statistics,
                then spreads throughout science and engineering. \\
\bottomrule
\end{tabular}
\end{center}

\section{Lessons for physics, computer science, and engineering}

\textbf{For physics:} The probability distribution was not an arbitrary curve.
It came from the canonical ensemble and encoded physical assumptions about
equilibrium, energy, and temperature. A simulation cannot be better than the
model and approximations on which it is based.

\textbf{For mathematics and statistics:} Stationarity and detailed balance
explain why the acceptance rule targets the desired distribution. Convergence,
autocorrelation, and sampling error determine whether a finite run provides a
useful estimate.

\textbf{For computer science:} The project required an algorithm whose local
operations were feasible on real hardware, a random-number mechanism, a careful
implementation, and data handling under severe memory constraints. Complexity
was changed by reformulating the task rather than merely running the original
calculation faster.

\textbf{For engineering:} MANIAC had to be designed, constructed, maintained,
and operated reliably. The method also required tradeoffs: proposal size versus
acceptance, physical realism versus tractability, and accuracy versus available
machine time.

\textbf{For everyone:} Modeling and simulation form a feedback loop. Physical
questions inspire mathematical structures; mathematical structures suggest
algorithms; hardware constraints reshape algorithms; computational results
expose weaknesses in models; and validation sends the team back around the loop.

\section{Questions for discussion}

\begin{enumerate}
  \item Which contribution was most essential: the physical question, Teller's
        ensemble idea, Rosenbluth's acceptance rule, Arianna Rosenbluth's
        program, or Metropolis's computer? Can this question have a unique answer?
  \item Why is a high-dimensional integral not solved simply by building a
        computer that is ten or one hundred times faster?
  \item In what sense is the Markov chain an artificial process rather than a
        simulation of the particles' real-time motion?
  \item What evidence would you require before trusting an equation of state
        computed by this method?
  \item Which forms of scientific labor are most likely to disappear from a
        method's name or from its standard textbook account?
  \item Identify a modern modeling problem in which advances in physics,
        algorithms, software, and hardware remain inseparable.
\end{enumerate}

\begin{thebibliography}{9}

\bibitem{metropolis1953}
N. Metropolis, A. W. Rosenbluth, M. N. Rosenbluth, A. H. Teller, and E. Teller,
``Equation of State Calculations by Fast Computing Machines,''
\emph{Journal of Chemical Physics} \textbf{21}, 1087--1092 (1953),
\href{https://doi.org/10.1063/1.1699114}{doi:10.1063/1.1699114}.

\bibitem{hastings1970}
W. K. Hastings,
``Monte Carlo Sampling Methods Using Markov Chains and Their Applications,''
\emph{Biometrika} \textbf{57}, 97--109 (1970),
\href{https://doi.org/10.1093/biomet/57.1.97}{doi:10.1093/biomet/57.1.97}.

\bibitem{gubernatis2005}
J. E. Gubernatis,
``Marshall Rosenbluth and the Metropolis Algorithm,''
\emph{Physics of Plasmas} \textbf{12}, 057303 (2005),
\href{https://doi.org/10.1063/1.1887186}{doi:10.1063/1.1887186}.

\bibitem{metropolis1987}
N. Metropolis,
``The Beginning of the Monte Carlo Method,''
\emph{Los Alamos Science} \textbf{15}, 125--130 (1987),
\url{https://mcnp.lanl.gov/pdf_files/Article_1987_LAS_Metropolis_125--130.pdf}.

\bibitem{anderson1986}
H. L. Anderson,
``Metropolis, Monte Carlo, and the MANIAC,''
\emph{Los Alamos Science} \textbf{14}, 96--108 (1986),
\url{https://mcnp.lanl.gov/pdf_files/Article_1986_LAS_Anderson_96--108.pdf}.

\bibitem{apsrosenbluth}
S. Chen,
``Arianna Rosenbluth and the Metropolis Monte Carlo Algorithm,''
\emph{APS News} (2022),
\url{https://www.aps.org/apsnews/2022/03/rosenbluth-metropolis-monte-carlo-algorithm}.

\bibitem{micibio}
Atomic Heritage Foundation,
``Augusta `Mici' Teller,''
\url{https://ahf.nuclearmuseum.org/ahf/profile/augusta-mici-teller/}.

\end{thebibliography}

\end{document}
